% ----------------------------------------------------------------------------
%                                   RESUMO
% ----------------------------------------------------------------------------
\begin{resumo}[RESUMO]

    \begin{SingleSpacing}
    
        O monitoramento contínuo de pacientes com a Doença de Parkinson exige sistemas computacionais eficientes para o suporte à decisão clínica. O sistema Painel Parkinson, originalmente desenvolvido como uma aplicação \textit{web}, auxilia profissionais da saúde por meio de \textit{dashboards} informativos. No entanto, a dependência exclusiva de navegadores \textit{web} e a centralização de processamento em nuvem geram limitações de desempenho, latência na renderização de dados complexos e elevados custos operacionais de hospedagem. Diante disso, este trabalho propõe a migração arquitetural do Painel Parkinson para uma aplicação \textit{desktop}. Utilizando o \textit{framework} Electron.js em conjunto com VueJS no \textit{front-end} e Python com PostgreSQL no \textit{back-end}, o projeto visa descentralizar o processamento de dados e otimizar a infraestrutura operacional. Metodologicamente, a investigação classifica-se como uma pesquisa aplicada, exploratória e descritiva, de abordagem qualitativa, fundamentada por meio de pesquisa bibliográfica para embasar o desenvolvimento do produto tecnológico voltado ao ambiente clínico. Espera-se que a nova versão \textit{desktop} entregue maior agilidade na interpretação de dados clínicos, melhore a experiência do usuário e viabilize a redução de gastos com servidores, garantindo um acompanhamento clínico mais acessível e preparado para futuras evoluções.
        
        \textbf{Palavras-chave}: Doença de Parkinson; Aplicação \textit{Desktop}; Otimização de Custos; Pesquisa Aplicada; Saúde Digital.
    
    \end{SingleSpacing}
    
\end{resumo}